%% ============================================================
\section{Architecture}
\label{sec:architecture}

SqC is implemented in Rust (2021 edition) and uses tree-sitter for
parsing.  The analysis pipeline proceeds in four stages, illustrated
in \cref{fig:pipeline}.

\begin{figure}[t]
\centering
\begin{tikzpicture}[
    node distance=0.6cm,
    box/.style={rectangle, draw, fill=blue!8, text width=5.8cm,
                minimum height=0.7cm, align=center, font=\footnotesize},
    arrow/.style={->, thick}
]
\node[box] (parse) {1. Tree-sitter Parse $\rightarrow$ Concrete Syntax Tree};
\node[box, below=of parse] (prescan) {2. Cross-file Prescan $\rightarrow$ Function Summaries};
\node[box, below=of prescan] (cfg) {3. Per-function CFG + Dataflow Analysis};
\node[box, below=of cfg] (rules) {4. Rule Evaluation (311 CERT C checkers)};
\draw[arrow] (parse) -- (prescan);
\draw[arrow] (prescan) -- (cfg);
\draw[arrow] (cfg) -- (rules);
\end{tikzpicture}
\caption{SqC analysis pipeline.  Each stage enriches the context
available to subsequent stages.}
\label{fig:pipeline}
\end{figure}

\subsection{Stage 1: Parsing with Tree-sitter}

SqC uses tree-sitter~\cite{tree-sitter} to produce a concrete syntax
tree (CST) for each C source file.  Unlike tools built on Clang's AST
(which requires a full compilation database), tree-sitter operates on
individual files without preprocessing.  This trades semantic precision
for practical deployability: SqC can analyze any C file without build
system integration.

The primary limitation is the inability to resolve preprocessor macros
and \texttt{\#include} directives.  SqC mitigates this through a
built-in database of approximately 370 standard library and platform
API functions (\cref{sec:stdlib-db}).

\subsection{Stage 2: Cross-File Prescan}
\label{sec:prescan}

When invoked with the \texttt{-d} (directory) flag, SqC performs a
pre-scan pass over all C files in the specified directories, collecting:

\begin{itemize}
    \item \textbf{Function definitions}: name, file, parameter count,
        and return type.
    \item \textbf{Function summaries}: per-parameter null state
        (DefinitelyNull, PossiblyNull, NotNull, Unknown) derived from
        call-site argument analysis across the project.
    \item \textbf{Local variable states}: within each function body,
        assignments to local variables are tracked to resolve argument
        values passed at call sites.
\end{itemize}

This was the single largest contributor to false positive reduction,
removing 148,622 FPs (31.2\%) in one iteration (\cref{sec:fp-cross-file}).

\subsection{Stage 3: Control-Flow Graphs and Dataflow}

For each function, SqC constructs a control-flow graph (CFG) handling
\texttt{if}/\texttt{else}, \texttt{switch}/\texttt{case}, loops with
\texttt{break}/\texttt{continue}, early returns, and \texttt{goto}
statements.

On the CFG, SqC performs forward dataflow analysis for null pointer
state tracking.  The null state lattice is:

\begin{equation}
\text{NullState} \in \{\text{Unknown}, \text{DefinitelyNull},
    \text{PossiblyNull}, \text{NotNull}\}
\end{equation}

Edge refinement propagates null information through branch conditions:
after \texttt{if (p != NULL)}, the true branch marks \texttt{p} as
NotNull while the false branch marks it as PossiblyNull.

The CFG infrastructure also supports:
\begin{itemize}
    \item A \textbf{constant evaluator} that resolves \texttt{sizeof}
        expressions, C standard limit macros (\texttt{INT\_MAX},
        \texttt{UINT64\_MAX}), file-scope \texttt{static const int}
        declarations, and compile-time arithmetic.
    \item \textbf{Dead-branch pruning}: when a branch condition
        resolves to a compile-time constant (literal, \texttt{\#define},
        or \texttt{static const}), only the feasible branch is included
        in the CFG, improving precision for all downstream analyses.
    \item \textbf{Initialization state tracking} for detecting
        use of uninitialized variables through conditional paths.
    \item \textbf{Value-range analysis} (VRA) for integer and
        floating-point overflow rules, tracking variable ranges through
        assignments, branches, and loop bounds.
\end{itemize}

\subsection{Stage 4: Rule Evaluation}

Each of the 311 rules is implemented as a Rust struct implementing
the \texttt{CertRule} trait.  Rules range from simple AST pattern
matchers (e.g., MSC24-C: obsolete function detection) to complex
multi-pass analyzers (e.g., EXP34-C: null dereference detection
using CFG-based dataflow with inter-procedural call-site propagation).

\textbf{UTF-8 source encoding.}  Tree-sitter represents all source
positions as byte offsets into the raw UTF-8 source file.  Industrial
C codebases routinely embed non-ASCII content in comments and string
literals: comments written in Chinese, Japanese, Arabic, Cyrillic, or
other scripts are common in embedded and international software
projects.  Rule implementations that inspect raw source text must
therefore use byte-indexed string operations throughout.
Character-counted indices produce incorrect slices---or runtime panics
in languages with strict Unicode safety---when a multi-byte sequence
falls before a slice boundary.  For example, a Chinese full-width
parenthesis (\texttt{U+FF08}, three UTF-8 bytes) encountered during
real-world analysis triggered a panic in the ARR00-C and PRE02-C
checkers, both of which were using \texttt{chars().enumerate()} indices
(character positions) as byte offsets.  SqC enforces the requirement
for byte-safe slicing through \texttt{char\_indices()}-based iteration
and includes UTF-8 regression test cases in each affected rule's test
suite.

\subsection{Standard Function Database}
\label{sec:stdlib-db}

Since tree-sitter cannot follow \texttt{\#include} directives, SqC
maintains a built-in database of approximately 370 known functions:
$\sim$180 C11 standard library, $\sim$90 POSIX, and $\sim$100 Windows
API functions.  This database serves multiple purposes: DCL31-C/DCL07-C
use it to avoid flagging standard library calls as undeclared; ERR33-C
uses it to identify functions whose return values must be checked;
EXP10-C uses a pure-function whitelist to exclude side-effect-free
calls from unsequenced evaluation warnings.
